\section{Estado del Arte}
\label{sec:estado_arte}

\subsection{Antecedentes sobre los videos de la FAC}
Los antecedentes más cercanos son dos trabajos desarrollados con videos FLIR suministrados por la Fuerza Aeroespacial Colombiana. Acosta-Bernal et al.~\cite{acosta2025eyes} extrajeron 1.657 fotogramas de un video de 777 segundos y entrenaron YOLOv11 para detectar edificaciones, vías, ríos, zonas de alteración del suelo y vehículos. Con filtrado de mediana y CLAHE, el recall pasó de 0,637 a 0,663 y el mAP@50 de 0,536 a 0,721. El trabajo reconoce que la cantidad limitada de escenas dificulta controlar la fuga de datos, y que SDZI fue la clase con peor desempeño (recall de 0,425 y mAP@50 de 0,402). De la Hoz~\cite{delahoz2025enhancing} mejoró EagleView AI, un prototipo surgido de un reto de innovación de la FAC, con una arquitectura modular para procesar videos extensos. En ambos, el preprocesamiento se usó como apoyo de la detección y no se evaluó de manera independiente.

\subsection{Limpieza de video aéreo}
Otros trabajos muestran que el preprocesamiento puede facilitar la detección en imágenes aéreas infrarrojas~\cite{jiang2023infrared,ozyurt2024contrast}, y que el ruido, el desenfoque y el bajo contraste afectan sobre todo a los objetos pequeños~\cite{hovhannisyan2025thermal,luo2025yolosmug,wang2026sgwdetr}.

Cuando no existen imágenes limpias de referencia se recurre a métodos autosupervisados. Noise2Noise demostró que una red puede entrenarse con pares de imágenes ruidosas independientes de la misma escena~\cite{lehtinen2018noise2noise}. Noise2Void evita la necesidad de pares mediante un punto ciego que oculta el píxel central~\cite{krull2019noise2void}, y StructN2V extiende la máscara para ruido con estructura espacial~\cite{broaddus2020structn2v}. Neighbor2Neighbor construye pares a partir de un único frame mediante submuestreo de píxeles vecinos~\cite{huang2021neighbor2neighbor} y Blind2Unblind aprovecha también los píxeles enmascarados~\cite{wang2022blind2unblind}. Para el ruido real de cámara, que está correlacionado entre píxeles vecinos, AP-BSN separa los píxeles antes de la red para romper esa correlación~\cite{lee2022apbsn} y MM-BSN combina varias máscaras~\cite{zhang2023mmbsn}. En video, Frame-to-Frame alinea frames contiguos con flujo óptico~\cite{ehret2019modelblind}, UDVD predice cada píxel a partir de sus vecinos espaciales y temporales~\cite{sheth2021udvd} y Frames2Residual desacopla la estructura de fondo del residuo dinámico~\cite{ji2026frames2residual}. Para eliminar información superpuesta, ProPainter combina propagación guiada por flujo y un transformer espacio-temporal~\cite{zhou2023propainter}.

En ausencia de referencia, la calidad se evalúa con métricas como NIQE~\cite{mittal2013making}, que mide la desviación respecto a estadísticas de imágenes naturales, y BRISQUE~\cite{mittal2012noreference}, basada en estadísticas locales de la imagen.

Los detectores localizan objetos de manera individual, lo que dificulta interpretar situaciones que dependen de la interacción entre elementos y de su evolución en el tiempo~\cite{cong2021spatialtemporal,wang2023context}. El \textit{dense video captioning} complementa esa detección al localizar segmentos temporales y describirlos en lenguaje natural~\cite{mkhallati2023soccernet,yang2023vid2seq,muksimova2025crossmodal}.

\subsection{Fuga de datos y particionamiento}
Cuando información del conjunto de prueba influye en el entrenamiento, ese conjunto deja de representar datos no vistos y la capacidad de generalización se sobreestima~\cite{sasse2025leakage}. Kapoor y Narayanan identifican ocho tipos de fuga y muestran que corregirlos puede cambiar las conclusiones de un estudio~\cite{kapoor2023leakage}. En sensores remotos, la correlación espacial entre muestras de entrenamiento y prueba sesga la estimación del error~\cite{hansch2017correct}.

Figueiredo y Mendes~\cite{figueiredo2024leakage} abordan este problema en conjuntos derivados de video: extraen características con CLIP, las reducen con t-SNE, comparan DBSCAN, OPTICS y agrupamiento aglomerativo, y asignan cada clúster completo a una única partición. Es el principal antecedente metodológico de este proyecto. Como representaciones visuales se consideran CLIP~\cite{radford2021clip} y DINOv2~\cite{oquab2024dinov2}; para la reducción de dimensionalidad, PaCMAP, que preserva relaciones locales y globales~\cite{wang2021pacmap}; y para el agrupamiento, HDBSCAN, que identifica grupos de formas irregulares y densidades variables~\cite{mcinnes2017hdbscan,glazner2025findleak}.

\subsection{Segmentación panóptica y supervisión puntual}
La mayoría de trabajos sobre minería y degradación ambiental tratan el problema como clasificación, detección o segmentación semántica de forma exclusiva, sobre imágenes satelitales~\cite{balaniuk2020mining,camalan2022change,deandrade2022deforestation,lemesneto2026sar}. La maquinaria también se ha detectado y segmentado desde plataformas aéreas~\cite{ersoz2024aidcon,liu2026lightweight}, pero de forma independiente de las regiones ambientales.

La segmentación panóptica~\cite{kirillov2019panoptic} unifica la segmentación semántica y la de instancias al representar a la vez categorías \textit{stuff} (regiones continuas) y \textit{things} (objetos individualizables), con Panoptic Quality como métrica conjunta. Panoptic-DeepLab~\cite{cheng2020panopticdeeplab} y Mask2Former~\cite{cheng2022mask2former} son referentes de esta formulación, y su traslado a imágenes aéreas enfrenta el costo de anotación, la escasez de datos y las variaciones de escala~\cite{decarvalho2022panoptic,cheng2024uavreview}.

Para reducir el costo de anotar a nivel de píxel se ha propuesto entrenar con anotaciones puntuales~\cite{wang2020weakly,wu2022dbfnet,chan2025sparse}, apoyándose en modelos como SAM para generar pseudo-etiquetas~\cite{metasam2}. Panoptic FCN~\cite{li2022panopticfcn} extiende la supervisión puntual a la segmentación panóptica y conserva buena parte del desempeño completamente supervisado con 20 puntos por objetivo; Point2Mask~\cite{li2023point2mask} la reduce a un único punto. Por último, incorporar las relaciones entre objetos y regiones mejora la interpretación de escenas aéreas~\cite{ren2024context} y, en particular, de escenas mineras con pocas muestras etiquetadas~\cite{dong2025mining}.

\subsection{Vacío de investigación}
Los métodos autosupervisados de reducción de ruido y de inpainting de video no se han integrado ni evaluado de manera sistemática sobre videos aéreos reales de vigilancia para la detección de minería ilegal; el particionamiento de estos datos no ha controlado la correlación entre fotogramas; y no se identificó un enfoque que combine segmentación panóptica, supervisión puntual y contexto espacial sobre imágenes FLIR de la Amazonía. Este proyecto aborda los tres vacíos en un mismo pipeline.
